\subsection{General Naming Requirements}

Names shall describe the purpose, behavior, or ownership of the object they identify. Abbreviations shall be avoided unless they are widely understood within the project.

Names shall:

\begin{itemize}
	\item use lowercase \texttt{snake\_case} for functions and variables;
	\item use uppercase \texttt{UPPER\_SNAKE\_CASE} for macros;
	\item use descriptive names for types;
	\item distinguish counts, sizes, indexes, capacities, and offsets;
	\item avoid names that differ only by capitalization;
	\item avoid names that shadow variables in an enclosing scope;
	\item avoid single-letter names except for short loop indexes;
	\item avoid misleading names.
\end{itemize}

\subsection{Functions}

Functions shall use verb-based names that describe their operation.

\begin{lstlisting}[language=C]
buffer_init();
buffer_write();
buffer_clear();
packet_is_valid();
\end{lstlisting}

Boolean-returning functions should use names such as:

\begin{lstlisting}[language=C]
is_valid();
has_error();
can_write();
\end{lstlisting}

\subsection{Variables}

Variables shall describe their content and, where useful, their unit.

\begin{lstlisting}[language=C]
size_t packet_length;
uint32_t timeout_ms;
uint8_t retry_count;
bool connection_is_open;
\end{lstlisting}

Names such as \texttt{data}, \texttt{value}, \texttt{temp}, and \texttt{result} should be avoided unless their meaning is clear from the local context.

\subsection{Types}

Project-defined types shall use a consistent project convention.

\begin{lstlisting}[language=C]
typedef struct buffer buffer_t;
typedef enum device_state device_state_t;
\end{lstlisting}

Structure tags and typedef names shall not unnecessarily duplicate one another.

\subsection{Constants and Macros}

Constants should preferably be represented using typed constants, enumerations, or \texttt{static const} objects instead of macros.

\begin{lstlisting}[language=C]
static const size_t max_packet_length = 1024U;
\end{lstlisting}

Macros shall use uppercase names with a project-specific prefix.

\begin{lstlisting}[language=C]
#define PROJECT_MAX_PACKET_LENGTH 1024U
\end{lstlisting}

\subsection{Private Names}

Private file-scope functions and variables shall use \texttt{static}. A project prefix may be used for private names when required to avoid collisions with external symbols.
